\documentclass[10pt]{article}
\usepackage[letterpaper, margin=0.6in]{geometry}
\usepackage{amsmath, amssymb}
\usepackage{booktabs}
\usepackage{xcolor}
\usepackage{titlesec}
\usepackage{enumitem}
\setlength{\parindent}{0pt}
\setlength{\parskip}{0.35em}
\renewcommand{\arraystretch}{1.15}
\titleformat{\section}{\large\bfseries}{\thesection}{0.5em}{}
\titlespacing*{\section}{0pt}{0.6em}{0.3em}
\pagestyle{empty}

\begin{document}

\begin{center}
{\Large\bfseries Q\&A cheat sheet --- skills vs.\ agents}\\
\smallskip
{\footnotesize Mercatus, 2026-06-09. \quad What's the difference, why both, and what's actually in this workflow.}
\end{center}

\hrule
\vspace{0.5em}

\section*{The one-sentence difference}

A \textbf{skill} is a Markdown checklist you trigger with a slash command, run by the Claude you're already talking to. An \textbf{agent} is a fresh Claude instance you spawn with its own focused context and persona, to which you delegate a self-contained task.

\section*{Skills}

\textbf{What:} a folder of plain-Markdown files at \texttt{.claude/skills/<name>/SKILL.md} that defines a multi-step procedure or checklist.\\
\textbf{How invoked:} the user types \texttt{/<skill-name>}.\\
\textbf{Where it runs:} in the main conversation, with the same Claude session you've been talking to.\\
\textbf{Best for:} repeatable workflows, codified protocols, quick shortcuts. ``How to compile this paper.'' ``How to validate the bibliography.'' ``How to commit, PR, and merge.''\\
\textbf{Cost:} tokens consume the main conversation's context window.

\section*{Agents}

\textbf{What:} a folder of Markdown files at \texttt{.claude/agents/<name>.md} that defines a sub-agent with a persona, a tool allowlist, and a system prompt.\\
\textbf{How invoked:} the main Claude calls the \texttt{Agent} tool, which spawns a new Claude instance with an isolated context.\\
\textbf{Where it runs:} a separate conversation. The sub-agent does its work, returns one final message, then disappears.\\
\textbf{Best for:} focused single-purpose review, deep judgement calls, parallel work, anything that would pollute the main context with raw output.\\
\textbf{Cost:} cheaper to the main context, but each agent run is its own Claude session.

\section*{Why you need both (four reasons)}

\begin{itemize}[leftmargin=1.4em, itemsep=0.15em]
  \item \textbf{Context hygiene.} A bibliography scan reads every \texttt{.tex} file in the manuscript. If that ran in your main session, your context would fill with raw bib entries. As an agent it stays clean.
  \item \textbf{Persona isolation.} The proofreader and the domain reviewer want different headspaces. Each agent gets its own system prompt and behaves like a different specialist. Skills don't change persona.
  \item \textbf{Parallelism.} You can spawn three agents on the same section simultaneously. Skills run one at a time in the main thread.
  \item \textbf{Workflow vs.\ judgement.} \texttt{/compile-latex} is three deterministic shell commands; that wants a skill. ``Did I represent Krugman accurately?'' is judgement; that wants an agent.
\end{itemize}

\section*{Agents currently in this workflow (4)}

\begin{tabular}{l p{0.74\textwidth}}
\toprule
Agent & What it does\\
\midrule
\texttt{proofreader}      & Grammar, typos, consistency, LaTeX hygiene visible from text. Read-only.\\
\texttt{domain-reviewer}  & Substantive review --- derivation correctness, assumption sufficiency, citation fidelity, code-theory alignment, logical consistency.\\
\texttt{r-reviewer}       & R code quality, reproducibility, figure-generation patterns, project theme compliance.\\
\texttt{verifier}         & End-to-end verification --- does the manuscript compile, do the R scripts run, are the output files non-empty.\\
\bottomrule
\end{tabular}

\section*{Skills currently in this workflow (12)}

\begin{tabular}{l p{0.72\textwidth}}
\toprule
Skill & What it does\\
\midrule
\multicolumn{2}{l}{\textit{Build / verify}}\\
\quad \texttt{/compile-latex}     & 3-pass XeLaTeX + bibtex on \texttt{Paper/main.tex}.\\
\quad \texttt{/validate-bib}      & Cross-reference \texttt{\textbackslash cite\{\}} keys against the .bib file.\\
\quad \texttt{/commit}            & Stage, commit, push, open PR, merge.\\
\multicolumn{2}{l}{\textit{Review (skills that invoke the agents above)}}\\
\quad \texttt{/proofread}         & Runs the proofreading protocol on the manuscript.\\
\quad \texttt{/review-paper}      & Comprehensive manuscript review --- argument, methods, citations, referee objections.\\
\quad \texttt{/review-r}          & R code review protocol on scripts.\\
\multicolumn{2}{l}{\textit{Research / writing workflows}}\\
\quad \texttt{/lit-review}        & Structured literature search and synthesis with citation extraction.\\
\quad \texttt{/research-ideation} & Generate testable hypotheses and empirical strategies from a topic.\\
\quad \texttt{/interview-me}      & Interactive interview to formalise a research idea into a spec.\\
\quad \texttt{/data-analysis}     & End-to-end R analysis from exploration through publication-ready tables.\\
\multicolumn{2}{l}{\textit{Project meta}}\\
\quad \texttt{/context-status}    & Show current context usage and session health.\\
\quad \texttt{/learn}             & Extract reusable knowledge from this session into a new skill.\\
\bottomrule
\end{tabular}

\section*{Two concrete examples of skills + agents working together}

\textbf{Example 1: gating a section before commit.}

\begin{enumerate}[leftmargin=1.5em, itemsep=0.05em, label=\arabic*.]
  \item Finish drafting \texttt{Paper/sections/09\_growth.tex}.
  \item In the main session: type \texttt{/review-paper}.
  \item The skill spawns the proofreader and domain-reviewer agents in parallel against the section.
  \item Each agent runs in its own context and returns a Markdown report to the main session.
  \item I read the two reports, decide which findings to act on, and edit accordingly.
  \item Type \texttt{/validate-bib} (skill, runs in main context --- it's just a regex sweep).
  \item Type \texttt{/compile-latex} (skill, runs the three shell commands).
  \item Type \texttt{/commit} (skill, stages and pushes).
\end{enumerate}
The two agents catch judgement-calls; the four skills handle deterministic workflow.

\medskip

\textbf{Example 2: when you'd use an agent vs.\ a skill for the same task.}

``Check the bibliography.''

\begin{itemize}[leftmargin=1.4em, itemsep=0.05em]
  \item \textbf{As a skill} (\texttt{/validate-bib}): regex extracts every \texttt{\textbackslash cite\{\}} key, compares to the .bib, reports missing/unused entries. Deterministic. Runs in 5 seconds in the main context.
  \item \textbf{As an agent} (substance-review): reads the .bib entries, checks whether each citation in the text accurately represents what the cited work actually argues. Judgement-heavy. Wants its own context.
\end{itemize}
Same surface task, two different jobs. Skill catches the mechanical gap (Smith1776\_wealth is in the text but not in the bib). Agent catches the substantive gap (you cited Krugman 1979 for a claim Krugman 1991 actually makes).

\vspace{0.4em}
\hrule
\vspace{0.4em}

\section*{Three verbal anchors}

\begin{itemize}[leftmargin=1.4em, itemsep=0.15em]
  \item \textit{Skill in one sentence:} ``A slash command that runs a checklist I wrote down, in the conversation I'm already in.''
  \item \textit{Agent in one sentence:} ``A fresh Claude with its own persona and context that I delegate a focused task to, and which sends one report back.''
  \item \textit{Why both:} ``Skills are deterministic workflow. Agents are judgement that wants its own headspace. The first stops me typing the same three commands every time. The second catches things a single Claude session would miss because it's juggling too much.''
\end{itemize}

\end{document}
